\documentclass[12pt]{article}
\usepackage{mc}
\mcheader{K--1}

% \pilerow{list of n}{box w}{box h}{gap}{per row}{sep}{radius}{label: 1 = numeral, 0 = none}{circle choices: 1 = yes}
\newcommand{\pilerow}[9]{%
  \begin{tikzpicture}
    \foreach \n [count=\j from 0] in {#1}{%
      \pgfmathsetmacro{\bx}{\j*(#2+#4)}%
      \draw[box] (\bx,0) rectangle ++(#2,-#3);
      \ifnum#8=1
        \counters{\bx+#2/2}{-#3/2+0.4}{\n}{#5}{#6}{#7}
        \node[font=\Large] at (\bx+#2/2,-#3+0.45) {\n};
      \else
        \counters{\bx+#2/2}{-#3/2}{\n}{#5}{#6}{#7}
      \fi
      \ifnum#9=1
        \node[font=\Large,anchor=base] at (\bx+#2*0.28,-#3-0.85) {1st};
        \node[font=\Large,anchor=base] at (\bx+#2*0.72,-#3-0.85) {2nd};
      \fi
    }
  \end{tikzpicture}%
}

% piles 1..12 in two rows of six
\newcommand{\twelve}{%
  \begin{tikzpicture}
    \foreach \n in {1,...,12}{%
      \pgfmathtruncatemacro{\col}{mod(\n-1,6)}%
      \pgfmathtruncatemacro{\row}{floor((\n-1)/6)}%
      \pgfmathsetmacro{\bx}{\col*3.0}%
      \pgfmathsetmacro{\by}{-\row*3.7}%
      \draw[box] (\bx,\by) rectangle ++(2.7,-3.3);
      \counters{\bx+1.35}{\by-1.25}{\n}{4}{0.56}{0.2}
      \node[font=\Large] at (\bx+1.35,\by-2.85) {\n};
    }
  \end{tikzpicture}%
}

% one row of n counters with gaps wide enough to draw rings: \ringrow{x}{y}{n}
\newcommand{\ringrow}[3]{%
  \foreach \i in {1,...,#3}{%
    \draw[ctr] ({#1+(\i-1)*1.35},#2) circle (0.28);
  }%
}

\begin{document}
\fontsize{15}{21}\selectfont

Play each game with a partner and a pile of counters. Take turns taking counters from the pile. Whoever takes the last counter wins.

\vspace{14pt}
\problem Take 1 or 2 counters on each turn. For each pile, circle whether you want to go 1st or 2nd.

\vspace{8pt}
\pilerow{2,3,4,5}{4.1}{4.0}{0.45}{4}{0.95}{0.36}{1}{1}

\vspace{40pt}
\problem It is your turn, and you may take 1 or 2 counters. Cross out the counters you would take to be sure to win.

\vspace{8pt}
\pilerow{4,5,7,8}{4.1}{4.4}{0.45}{4}{0.95}{0.38}{1}{0}

\newpage
\problem Take 1 or 2 counters on each turn. Color every pile where you would rather go second.

\vspace{8pt}
\twelve

\vspace{40pt}
\problem Ben always takes 2 counters when he can, and you take 1 or 2. Circle each pile where you can beat Ben when he goes first.

\vspace{8pt}
\pilerow{2,3,4,5,6,7}{2.7}{3.6}{0.3}{4}{0.58}{0.21}{1}{0}

\newpage
\problem Now take 1, 2, or 3 counters on each turn. Color every pile where you would rather go second.

\vspace{8pt}
\twelve

\vspace{40pt}
\problem Take 1 or 3 counters on each turn. Color every pile where you would rather go second.

\vspace{8pt}
\twelve

\newpage
\problem Start with a pile of 20 counters. For each rule, circle whether you want to go 1st or 2nd.

\vspace{8pt}
\begin{tikzpicture}
  \draw[box] (0,0) rectangle (5.4,-5.0);
  \counters{2.7}{-2.1}{20}{5}{0.9}{0.32}
  \node[font=\Large] at (2.7,-4.5) {20};
  \foreach \r/\t in {0/{Take 1 or 2.},1/{Take 1, 2, or 3.},2/{Take 1 or 3.}}{%
    \node[anchor=base west,font=\Large] at (6.8,-1.0-\r*1.7) {\t};
    \node[anchor=base,font=\Large] at (12.8,-1.0-\r*1.7) {1st};
    \node[anchor=base,font=\Large] at (14.8,-1.0-\r*1.7) {2nd};
  }
\end{tikzpicture}

\vspace{40pt}
\problem Take 1 or 2 counters from one of the two piles on each turn. For each pair of piles, circle whether you want to go 1st or 2nd.

\vspace{8pt}
\begin{tikzpicture}
  \foreach \a/\b [count=\j from 0] in {2/2,1/2,3/3,1/4}{%
    \pgfmathsetmacro{\bx}{\j*4.55}%
    \draw[box] (\bx,0) rectangle ++(4.1,-3.8);
    \draw[line width=0.6pt, black!50] (\bx+2.05,-0.3) -- (\bx+2.05,-3.5);
    \counters{\bx+1.025}{-1.5}{\a}{2}{0.75}{0.28}
    \counters{\bx+3.075}{-1.5}{\b}{2}{0.75}{0.28}
    \node[font=\Large] at (\bx+1.025,-3.25) {\a};
    \node[font=\Large] at (\bx+3.075,-3.25) {\b};
    \node[font=\Large,anchor=base] at (\bx+4.1*0.28,-4.65) {1st};
    \node[font=\Large,anchor=base] at (\bx+4.1*0.72,-4.65) {2nd};
  }
\end{tikzpicture}

\newpage
\noindent
\begin{minipage}[t]{0.47\textwidth}\raggedright
\problem A game starts with 4 counters, and each turn takes 1 or 2. Draw rings around the counters taken on each turn to show every different way the game can go.
\end{minipage}\hfill
\begin{minipage}[t]{0.47\textwidth}\raggedright
\problem Now find every different way a game with 5 counters can go, taking 1 or 2 on each turn.
\end{minipage}

\vspace{12pt}
\noindent
\begin{minipage}[t]{0.47\textwidth}
\begin{tikzpicture}[baseline=(current bounding box.north)]
  \draw[box] (0,0) rectangle (6.0,-11.0);
  \foreach \r in {0,...,6}{%
    \ringrow{1.05}{-0.8-\r*1.55}{4}
  }
\end{tikzpicture}
\end{minipage}\hfill
\begin{minipage}[t]{0.47\textwidth}
\begin{tikzpicture}[baseline=(current bounding box.north)]
  \draw[box] (0,0) rectangle (7.6,-15.65);
  \foreach \r in {0,...,9}{%
    \ringrow{1.1}{-0.8-\r*1.55}{5}
  }
\end{tikzpicture}
\end{minipage}

\end{document}
